\chapter{ساخت ربات با بات‌فادر}

\begin{goalbox}
حساب ربات تلگرام را بسازید و توکن آن را به‌شکل امن در اختیار پایتون بگذارید.
\end{goalbox}

\section{درخواست ربات از بات‌فادر}

\begin{enumerate}
\item در تلگرام حساب تأییدشدهٔ
\lr{@BotFather}
را جست‌وجو کنید.
\item گفت‌وگو را آغاز کنید و دستور
\lr{/newbot}
را بفرستید.
\item یک نام نمایشی وارد کنید؛ برای نمونه
\lr{Alireza Tutorial Bot}
مناسب است.
\item نام کاربری یکتایی وارد کنید که به
\lr{bot}
پایان برسد؛ برای نمونه
\lr{alireza\_tutorial\_123\_bot}
را امتحان کنید. اگر قبلاً گرفته شده است، نام دیگری بسازید.
\item بات‌فادر یک توکن می‌دهد. با آن مانند رمز عبور رفتار کنید.
\end{enumerate}

\begin{warningbox}
توکن واقعی را در کد، تصویر، پیام یا گیت قرار ندهید. اگر توکن فاش شد، دستور
\lr{/revoke}
را برای بات‌فادر بفرستید و توکن تازه‌ای بسازید.
\end{warningbox}

\section{تنظیم توکن برای یک نشست ترمینال}

مقدار نمونه را با توکنی که از بات‌فادر گرفته‌اید جایگزین کنید.

در مک یا لینوکس:
\begin{LTR}
\begin{lstlisting}[language=bash,numbers=none]
export BOT_TOKEN="paste_your_token_here"
\end{lstlisting}
\end{LTR}

در پاورشل ویندوز:
\begin{LTR}
\begin{lstlisting}[language={},numbers=none]
$env:BOT_TOKEN="paste_your_token_here"
\end{lstlisting}
\end{LTR}

با بسته‌شدن آن ترمینال، متغیر حذف می‌شود. این روش برای نخستین پروژه ساده و امن
است. در هر نشست تازه باید آن را دوباره تنظیم کنید.

\begin{checkpointbox}
فقط وجود متغیر را بررسی کنید و مقدار محرمانهٔ آن را نمایش ندهید.
\begin{LTR}
\begin{lstlisting}[language=bash,numbers=none]
python -c "import os; print('token set' if os.getenv('BOT_TOKEN') else 'missing')"
\end{lstlisting}
\end{LTR}
نتیجهٔ مورد انتظار عبارت
\lr{token set}
است.
\end{checkpointbox}
